\chapter{计算机体系结构与组织}

\section{领域导言与演进背景}

\begin{preamblebox}[知识领域概述与变革]
在计算机学科的工程实践中，专业人员通常花费大量时间编写高能效代码，以解决特定应用领域的复杂计算问题。随着摩尔定律（Moore's Law）与登纳德缩放定律面临物理极限与终结，几乎所有计算系统（包括最普及的商用消费级硬件）均日益依赖硬件系统层面的并行性来满足持续增长的计算性能需求。这种从传统的纯顺序处理向高度并行计算的根本范式转变，要求软件从业者对底层计算机体系结构建立更为深刻透彻的理解。体系结构已不再能被视作一个可以任意套用抽象原则的“黑盒”；相反，现代程序员必须能够“透视黑盒内部”，主动利用特定的底层硬件架构组件与微架构特性，以系统性提升计算性能并大幅优化能源效率。

“计算机体系结构与组织”（Architecture and Organization, AR）知识领域旨在培养学生对现代一切计算基石——底层硬件运行环境，以及硬件向高层软件栈所提供各类接口的深层次理解。本领域涵盖的目标硬件跨度广泛，从低功耗物联网嵌入式处理器，直至面向超算与云端的大规模企业级多处理器系统。

深入学习本领域的知识主题，将使学生全面领会现代计算机系统的核心体系结构设计原则，特别是如何攻克驾驭硬件并行性以维持未来持续性能与能效提升的重大挑战。此外，深刻理解计算机体系结构是开发兼具高性能与超低能耗程序的必由之路。本知识领域将协助计算机科学专业学生彻底告别“黑盒化”思维，建立对底层计算机系统的全局敏锐洞察，并充分发掘与释放特定硬件架构所能实现的极致效能。

\textbf{自 CS2013 以来的核心变革：}
\begin{enumerate}
    \item \textbf{核心知识单元全面重构与更新}：重点重塑了“体系结构/存储体系”（体系结构/层次存储系统）与“体系结构/性能与能效”（体系结构/性能与能效）两个知识单元，深度融合了片上多级缓存优化、高带宽存储以及能耗感知计算的最新前沿进展。
    \item \textbf{新增“异构架构”核心单元（体系结构/异构架构）}：直面计算机体系结构领域的新兴突破，系统覆盖存算一体 (PIM)、领域专用架构（Domain-Specific Architectures, DSA，如神经网络处理器 NPU/TPU）、片上互联与先进封装技术（3DIC / Chiplets 芯粒）等前沿架构。
    \item \textbf{新增“量子体系结构”先锋单元（体系结构/量子体系结构）}：鉴于量子计算与量子处理器架构研究的飞速推进，专门增设该单元作为前沿“工具箱”，系统涵盖量子比特物理实现、量子线路、量子门、基础量子算法及量子处理单元（QPU）的微架构组织。
    \item \textbf{整合交叉重叠单元}：将原 CS2013 中的“多处理与替代架构”（体系结构/Multiprocessing and Alternative Architectures）整合并入全新的“异构架构”单元中，使知识逻辑更加紧凑清晰。
\end{enumerate}
\end{preamblebox}

\section{核心学时分配建议}

本知识领域在“向日葵模型”下的学时推荐如下表所示：

\begin{table}[H]
\centering
\setlength{\tabcolsep}{3.5pt}\fontsize{8pt}{9.6pt}\selectfont\begin{tabularx}{\columnwidth}{Y c c}
\toprule
\textbf{知识单元} & \textbf{计算机核心} & \textbf{领域核心} \\
\midrule
数字逻辑与数字系统 & -- & 3 \\
机器级数据表示 & 1 & -- \\
汇编级机器组织 & 1 & 2 \\
层次存储系统 & 6 & -- \\
接口与通信 & 1 & -- \\
功能部件组织 & -- & 2 \\
性能与能效 & -- & 3 \\
异构架构 & -- & 3 \\
量子体系结构 & -- & 3 \\
\midrule
\textbf{合计} & \textbf{9 学时} & \textbf{16 学时} \\
\bottomrule
\end{tabularx}
\end{table}

\section{知识单元与详细大纲}

\subsection{数字逻辑与数字系统}
\textbf{【学时建议】}：3 领域核心学时。

本单元建立现代计算系统的底层数字电路与硬件抽象基础：

\subsubsection{领域核心主题}
\begin{itemize}
    \item \textbf{体系结构历史与演进概述}：计算技术的发展历程（电子管、晶体管、集成电路至超大规模集成电路 VLSI），冯·诺依曼结构的发展演变。
    \item \textbf{组合逻辑与时序逻辑 / 现场可编程门阵列 (FPGA)}：
    \begin{itemize}
        \item 基础组合逻辑构建块：逻辑门、多路选择器、译码器、优先编码器与算术逻辑单元（ALU）。
        \item 时序逻辑基本构建块：双稳态锁存器、触发器（Flip-Flop，D 触发器与 JK 触发器）、时序寄存器、计数器与有限状态机（FSM）。
    \end{itemize}
    \item \textbf{多层硬件与软件抽象体系结构}：晶体管层、门级电路、寄存器传输级（RTL）、指令集体系结构（ISA）、操作系统内核及应用程序层。
    \item \textbf{计算机辅助设计 (CAD) 工具链}：处理硬件表示与体系结构模拟的综合、布局布线与仿真工具。
    \item \textbf{高层次综合}：
    \begin{itemize}
        \item 寄存器传输级表示法。
        \item 硬件描述语言（HDL，如 Verilog、VHDL 与现代化 Chisel 语言）。
    \end{itemize}
    \item \textbf{片上系统 设计流程}：软硬件协同设计、IP 核复用、时钟树综合与物理验证流程。
    \item \textbf{物理与电气工程约束}：
    \begin{itemize}
        \item 门延迟与建立/保持时间。
        \item 扇入与扇出能力。
        \item 能量与功耗约束（动态开关功耗与静态漏电功耗）。
        \item 信号传播物理极限（光速在硅介质中的传播延迟）。
    \end{itemize}
\end{itemize}

\subsubsection{示范性学习成果与能力目标}
\begin{enumerate}
    \item \textbf{[解释]} 评述计算机底层组件技术的历史演进——从真空电子管到现代超大规模集成电路（VLSI），从传统单体大型机架构演化至现代仓储级计算机 (WSC)的分布式数据中心集群架构。
    \item \textbf{[评估]} 深入分析现代异构计算机体系结构中组件之间及组件内部的并行性形态与数据依赖约束。
    \item \textbf{[解释]} 详尽阐释“功耗墙”（Power Wall）产生的物理根源，剖析其为何终结了单核频率的提升，并成为全面发掘与驾驭硬件并行性的核心瓶颈。
    \item \textbf{[开发]} 提出计算机基础构建模块的严密设计方案：算术逻辑单元 ALU（门级实现）、寄存器与寄存器堆（门级实现）、中央处理器 CPU（寄存器传输级 RTL 实现）以及各级存储器（寄存器传输级 RTL 实现）。
    \item \textbf{[评估]} 评估并验证简易计算机硬件设计的基础构建模块（如算术逻辑单元、寄存器组及内部总线数据流动）的时序与逻辑正确性。
    \item \textbf{[评估]} 验证简易处理器的时序图行为特征，精准识别并定位微架构中的数据依赖与时序冒险风险。
\end{enumerate}

\subsection{机器级数据表示}
\textbf{【学时建议】}：1 计算机核心学时。

本单元建立机器底层二进制数据存储、编码与精度约束的数理理解：

\subsubsection{计算机核心主题}
\begin{itemize}
    \item \textbf{位、字节与字}：二进制位、字节与机器字长的体系结构定义，自然对齐与非对齐访存。
    \item \textbf{数值数据表示与进位制系统}：
    \begin{itemize}
        \item 定点数表示法。
        \item 浮点数表示法（Floating-point representation，IEEE 754 标准：阶码、尾数、规格化数、非规格化数、NaN 与无穷大值）。
    \end{itemize}
    \item \textbf{有符号数与补码表示}：原码、反码与补码（Two's-complement）表示法，算术溢出判定与符号位扩展机制。
    \item \textbf{非数值数据的位级表示}：ASCII 与 Unicode（UTF-8、UTF-16）字符集编码规范，图像、音频等多媒体数据的位级量化表征。
    \item \textbf{复合数据结构的内存布局}：结构体/记录的字节对齐与填充规则，多维数组在内存中的行优先/列优先连续排布。
\end{itemize}

\subsubsection{示范性学习成果与能力目标}
\begin{enumerate}
    \item \textbf{[解释]} 深刻阐释现代计算机体系中“一切皆为数据”（包括程序指令自身亦以二进制数据形态存储执行）的冯·诺依曼同构存储哲学及其技术影响。
    \item \textbf{[应用]} 结合数值表示位分布图，标定与剖析由于定长二进制表示局限而引发精度损失、舍入误差、数值截断及上下溢出的敏感区间。
    \item \textbf{[解释]} 评述原码与补码在存储负整数时的底层机制差异，阐明补码在统一加减法硬件运算与消除双零歧义上的计算优势。
    \item \textbf{[解释]} 结合实际工程场景，合理论证并解释不同数值精度格式（如 FP64、FP32、FP16、BF16 及 INT8）在科学计算与深度学习模型推理中的选型依据。
    \item \textbf{[解释]} 阐述非数值数据（如字符、字符串、结构体记录与多维数组）在底层内存中的位级排布、对齐填充与指针寻址机制。
    \item \textbf{[应用]} 熟练完成不同数制（二进制、八进制、十进制、十六进制）与数值表示格式之间的精确手工与程序化相互转换。
\end{enumerate}

\subsection{汇编级机器组织}
\textbf{【学时建议】}：1 计算机核心学时，2 领域核心学时。

本单元建立从高级语言到机器底层指令执行与微架构映射的核心桥梁：

\subsubsection{计算机核心主题}
\begin{itemize}
    \item \textbf{冯·诺依曼计算机体系结构}：存储程序概念、程序计数器（PC）、总线瓶颈与核心功能子系统（计算、控制、存储、输入、输出）划分。
    \item \textbf{控制单元与指令周期}：取指 (IF)、译码 (ID)与执行 (EX)的微周期节拍机制。
    \item \textbf{费林分类法 (Flynn's Taxonomy) 导论}：单指令流单数据流（SISD）、单指令流多数据流（SIMD）与多指令流多数据流（MIMD）的体系结构特征与演进。
    \item \textbf{共享内存多处理器 / 多核组织架构}：多核对称多处理器（SMP）、片上互联与共享总线微架构（与操作系统 OS 交叉）。
\end{itemize}

\subsubsection{领域核心主题}
\begin{itemize}
    \item \textbf{指令集体系结构 (ISA)}（涵盖典型 CISC 与 RISC 架构代表：x86、ARM 与开源开放的 RISC-V）：
    \begin{itemize}
        \item 指令编码格式（操作码 Opcode、寄存器操作数、立即数）。
        \item 数据传送、算术与逻辑运算、程序控制流指令及基本 I/O 指令。
        \item 寻址模式（直接寻址、立即寻址、间接寻址、基址变址寻址与 PC 相对寻址）。
        \item 机器语言编程与机器码格式。
        \item 汇编语言程序设计技术。
    \end{itemize}
    \item \textbf{子程序调用与返回机制}：调用约定、参数传递、返回值保存、栈帧平衡与寄存器保护机制（调用者保存 vs 被调用者保存，交叉引用 PL/编程语言翻译与执行，与 OS 交叉）。
    \item \textbf{I/O 与中断处理机制}：中断请求（IRQ）、中断向量表（IVT）与保护现场（与 OS 交叉）。
    \item \textbf{程序运行时内存段划分}：堆区、静态数据区、栈区与代码段（与 OS 交叉）。
\end{itemize}

\subsubsection{示范性学习成果与能力目标}
\begin{enumerate}
    \item \textbf{[解释]} 在现代嵌入式系统语境下（尤其是紧耦合片上内存与高延迟片外存储系统）分析并定位经典冯·诺依曼功能单元的物理映射与作用。
    \item \textbf{[解释]} 详细阐述经典冯·诺依曼架构下指令执行的微周期阶段，以及针对多线程并发、多处理器硬件同步机制与 SIMD 向量化执行的体系结构拓展支持。
    \item \textbf{[评估]} 在包含指令级并行（ILP）与流水线冒险（结构冒险、数据冒险与控制冒险）的示例微架构图上进行专业标注，评述现代处理器流水线化解冒险的关键技术手段。
    \item \textbf{[解释]} 阐明机器指令在二进制机器码层级以及符号汇编器语境下的表示与转换映射规则。
    \item \textbf{[应用]} 将常见的高级程序语言控制结构（如 if-else 条件分支、for/while 循环、函数调用与对象成员访问）映射并翻译为等价的汇编/机器指令序列。
    \item \textbf{[评估]} 评述与对比不同指令格式的设计权衡，例如单指令操作数地址数（零地址、一地址、二地址、三地址架构）以及定长指令集（RISC）与变长指令集（CISC）的优劣。
    \item \textbf{[解释]} 结合栈帧剖析图，详细阐述汇编语言层面子程序调用的压栈、参数传递、局部变量分配、返回地址维护与现场恢复的全过程。
    \item \textbf{[解释]} 深刻阐释硬件中断机制、中断向量表（IVT）与 I/O 端口操作的基本工作原理。
    \item \textbf{[开发]} 编写一段高效的汇编语言程序，实现严密底层的字符串遍历、模式查找与文本内存操作。
\end{enumerate}

\subsection{层次存储系统}
\textbf{【学时建议】}：6 计算机核心学时。

本单元是计算机体系结构的核心支柱，深入解析现代多级存储系统如何化解“访存墙”瓶颈：

\subsubsection{计算机核心主题}
\begin{itemize}
    \item \textbf{存储体系原理}：层次存储金字塔，时间局部性与空间局部性的决定性意义（与 OS 交叉）。
    \item \textbf{主存储器组织与操作}：动态随机存取存储器（DRAM）微架构、电容电荷刷新机制、行缓冲、双倍数据速率（DDR）技术与多通道总线配置。
    \item \textbf{持久化与非易失性存储}：固态硬盘（SSD）、NAND Flash 闪存存储原理、磨损均衡与传统机械磁盘（HDD）组织（与 OS 交叉）。
    \item \textbf{存储性能度量}：访存延迟、周期时间、峰值与持续带宽以及多体交叉存取。
    \item \textbf{高速缓存 架构原理}（与 OS 交叉）：
    \begin{itemize}
        \item 地址映射策略：直接映射、全相联与组相联。
        \item 缓存块大小的性能权衡与末级缓存（LLC）设计。
        \item 替换策略：最近最少使用（LRU）、先进先出（FIFO）、伪 LRU（PLRU）与随机替换。
        \item 写策略：写直达与写回，写分配与非写分配。
    \end{itemize}
    \item \textbf{多处理器缓存一致性}：假共享（False Sharing）、总线监听协议、基于目录的协议与 MESI / MOESI 状态机模型。
    \item \textbf{虚拟内存 体系结构支持}：页表结构、多级页表、转译后备缓冲器（TLB 快表）、缺页异常处理流程与内存保护机制（与 OS 交叉）。
    \item \textbf{硬件故障处理与可靠性保证}（与 OS 交叉）。
    \item \textbf{硬件级可靠性技术}（交叉引用 系统基础/基于冗余的可靠性）：
    \begin{itemize}
        \item 差错控制编码：奇偶校验位、ECC 内存校验、海明码与单错校正双错检测（SEC-DED）。
        \item 硬件辅助数据压缩。
        \item 端到端数据完整性验证。
    \end{itemize}
    \item \textbf{非冯·诺依曼架构探索}：
    \begin{itemize}
        \item 存算一体 / 内存计算与近存计算。
    \end{itemize}
\end{itemize}

\subsubsection{示范性学习成果与能力目标}
\begin{enumerate}
    \item \textbf{[解释]} 借助计算机存储层次金字塔图，辨析主流半导体与非易失存储技术（如片上寄存器、SRAM、DRAM、NAND Flash、持久内存）的物理特性、相对每比特成本、容量密度与访存延迟差异。
    \item \textbf{[评估]} 精确测量与量化不同层级存储访问延迟（从 L1 Cache 到主存及持久磁盘）对真实应用软件端到端执行时间的影响。
    \item \textbf{[解释]} 系统阐明具备硬件虚拟内存管理单元（MMU）的现代系统所提供的关键功能（地址空间隔离、动态内存保护、虚拟空间超额分配及按需分页）。
    \item \textbf{[应用]} 熟练推导并计算在多样化 Cache 与内存微架构配置、不同指令与数据引用混合比场景下的平均访存时间。
\end{enumerate}

\subsection{接口与通信}
\textbf{【学时建议】}：1 计算机核心学时。

本单元探讨处理器与外设之间的高速互联、数据交互与通信控制：

\subsubsection{计算机核心主题}
\begin{itemize}
    \item \textbf{输入/输出 (I/O) 基础原理}（与 OS 及 SPD 交叉）：
    \begin{itemize}
        \item 异步握手协议与多级 FIFO 缓冲机制。
        \item 程序控制 I/O（Programmed I/O / 轮询 Polling）。
        \item 中断驱动 I/O（Interrupt-driven I/O）。
    \end{itemize}
    \item \textbf{中断系统微架构}：向量中断与优先级中断仲裁、中断确认握手与多重中断嵌套（与 OS 交叉）。
    \item \textbf{外存储物理组织与通信驱动}：磁盘驱动器物理结构、固态存储控制器通信协议（如 NVMe、SATA）（与 OS 交叉）。
    \item \textbf{系统总线基础技术}（与 OS 及 SPD 交叉）：
    \begin{itemize}
        \item 总线时序协议：同步总线与异步总线。
        \item 总线仲裁机制：集中式仲裁（菊花链、独立请求）与分布式仲裁。
        \item 直接内存访问（DMA: Direct-Memory Access）：DMA 控制器微架构、突发传输模式与周期窃取。
    \end{itemize}
    \item \textbf{片上网络}：多核片上互联拓扑结构（2D Mesh、Torus、Crossbar）、虚通道流控机制与微路由器仲裁。
\end{itemize}

\subsubsection{示范性学习成果与能力目标}
\begin{enumerate}
    \item \textbf{[解释]} 参照中断控制微架构时序图，系统剖析硬件中断如何协调 I/O 外设控制器与 CPU 之间的数据传送与突发事件响应。
    \item \textbf{[解释]} 详细列举现代计算机系统中的各类互联总线与接口标准（如 PCIe、USB、AXI 片上总线），对比其时钟频率、通道宽度与有效带宽特性。
    \item \textbf{[评估]} 系统梳理传统机械硬盘（HDD）的优劣特征，并从寻道延迟、随机读写 IOPS、抗震性、耐用磨损及单位存储成本等维度与固态硬盘（SSD）进行深度对比评估。
\end{enumerate}

\subsection{功能部件组织}
\textbf{【学时建议】}：2 领域核心学时。

本单元深入剖析处理器内部数据通路构建、流水线调度与控制单元设计：

\subsubsection{领域核心主题}
\begin{itemize}
    \item \textbf{基础数据通路 的设计与实现}：ALU 算术逻辑单元、寄存器堆内部总线互联、时钟周期控制、单周期与多周期数据通路划分。
    \item \textbf{控制单元 的设计与实现}：
    \begin{itemize}
        \item 硬布线逻辑实现：组合逻辑门阵列、有限状态机时序生成器。
        \item 微程序控制器实现：控制存储器、微指令格式（水平型微指令与垂直型微指令）、微地址转移逻辑与微解释器。
    \end{itemize}
    \item \textbf{指令流水线技术}：经典 RISC 五级流水线（取指 IF、译码 ID、执行 EX、访存 MEM、写回 WB）及其时空图。
    \item \textbf{流水线冒险及其硬件消解}：
    \begin{itemize}
        \item 结构冒险：硬件资源冲突与多端口存储器解决方案。
        \item 数据冒险：写后读（RAW）、读后写（WAR）、写后写（WAW），旁路转发机制与硬件流水线暂停（Stalling / 流水线气泡 Bubble）。
        \item 控制冒险：分支跳转延迟、分支延迟槽与静态分支预测。
    \end{itemize}
    \item \textbf{指令级并行 导论}：多发射流水线（Superscalar 超标量架构基础）。
\end{itemize}

\subsubsection{示范性学习成果与能力目标}
\begin{enumerate}
    \item \textbf{[评估]} 验证并评述现代计算机微架构中单周期、多周期与经典多级流水线数据通路实现的优缺点与关键时钟频率约束。
    \item \textbf{[开发]} 分别采用硬布线控制器组合逻辑与微程序微指令控制字方案，针对两整型操作数相加并写回寄存器的微操作序列，设计出一组精确完整的硬件控制信号集合。
    \item \textbf{[解释]} 深刻阐释利用流水线发掘指令级并行的基本机理，系统分析运行过程中可能产生的各类严重流水线冒险及其硬件规避对策。
    \item \textbf{[开发]} 独立设计一个包含完整数据通路与控制单元、支持基础 RISC 指令子集的简易可综合处理器微架构。
    \item \textbf{[评估]} 针对给定的处理器流水线架构与层次内存延迟参数，精确推导并计算程序的平均每条指令周期数（CPI）及处理器的实际指令执行吞吐量。
\end{enumerate}

\subsection{性能与能效}
\textbf{【学时建议】}：3 领域核心学时。

本单元深入研究微架构前沿性能挖掘技术、能效优化策略与芯片物理约束：

\subsubsection{领域核心主题}
\begin{itemize}
    \item \textbf{性能与能耗综合评估导论}：计算性能三要素公式（$T = \text{IC} \times \text{CPI} \times \tau$）、动态功耗（$P_{\text{dynamic}} \propto C V^2 f$）与静态漏电功耗、内存访问能耗与片间通信开销。
    \item \textbf{高级微架构推测与乱序执行}：
    \begin{itemize}
        \item 分支预测技术：静态预测、动态 2 位饱和计数器、局部/全局自适应分支预测器与锦标赛预测器。
        \item 推测执行与精确异常处理。
        \item 乱序执行（Out-of-Order Execution）：Tomasulo 算法、保留站、重排序缓冲区 (ROB)与寄存器重命名。
    \end{itemize}
    \item \textbf{硬件预取技术}：指令预取、基于步长的数据预取器。
    \item \textbf{向量处理器与 GPU 架构增强}：单指令多数据（SIMD）向量扩展指令集、单指令多线程（SIMT）模型、张量加速核心。
    \item \textbf{多线程的硬件支持}：
    \begin{itemize}
        \item 细粒度、粗粒度多线程与同步多线程（SMT / 超线程技术）。
        \item 硬件级竞态条件防范。
        \item 硬件锁机制的实现：原子比较并交换（CAS）、加载链接/条件存储（LL/SC）与硬件自旋锁。
        \item 点对点硬件同步与硬件屏障同步机制。
    \end{itemize}
    \item \textbf{体系结构可扩展性定律}：阿姆达尔定律（Amdahl's Law）与古斯塔夫森定律（Gustafson's Law）。
    \item \textbf{非传统替代架构}：超长指令字（VLIW / EPIC 显式并行指令计算）、专用硬件加速器与协处理器。
    \item \textbf{动态电压与频率调节 (DVFS: Dynamic Voltage and Frequency Scaling)}。
    \item \textbf{“暗硅”挑战}：登纳德缩放定律失效后芯片热耗散极限（TDP）与不能全速同时点亮所有晶体管的物理困境。
\end{itemize}

\subsubsection{示范性学习成果与能力目标}
\begin{enumerate}
    \item \textbf{[解释]} 深刻评述衡量现代计算机系统能效与计算性能的核心基准测试与度量指标（如 SPEC 基准套件、MIPS/MFLOPS、每瓦性能比 Performance-per-Watt、能耗延迟积 EDP）。
    \item \textbf{[评估]} 跟踪推测执行与乱序微架构状态转移图，系统撰写并剖析在分支预测正确提交与预测失配重滚刷新的关键判定决策。
    \item \textbf{[开发]} 搭建并绘制针对特定 GPU 硬件架构的性能-功耗基准测试测评体系图。
    \item \textbf{[开发]} 编写一个底层多线程并发程序（如基于 C++ / Python），实现两高维整型向量的并行高效相加，并分析 CPU 核心绑定与访存开销。
    \item \textbf{[开发]} 针对特定应用计算场景（如边缘端低功耗推理或高并发网络包解析），提出一组关于替代性硬件架构的权衡选型方案。
    \item \textbf{[解释]} 深入阐释动态电压与频率调节（DVFS）的核心物理原理、功耗立方/平方降比效应及在现代操作系统调度器中的应用机制。
    \item \textbf{[评估]} 精确评测并量化神经网络模型在采用 8 位整型（INT8）量化相较于传统 32 位浮点（FP32）格式时所获得的内存带宽压缩率与硬件能耗节约幅度。
\end{enumerate}

\subsection{异构架构}
\textbf{【学时建议】}：3 领域核心学时。

本单元面向后摩尔时代前沿，系统解构多类异构加速器与先进微系统集成技术：

\subsubsection{领域核心主题}
\begin{itemize}
    \item \textbf{异构并行体系结构形态}：SIMD 向量处理器、通用图形处理器（GPGPU）、张量处理器（TPU）与神经网络处理器（NPU）微架构解构。
    \item \textbf{异构存储体系结构}：
    \begin{itemize}
        \item 共享内存与分布式内存混合拓扑。
        \item 易失性与非易失性混合持久存储。
        \item 异构缓存一致性协议（CXL, Compute Express Link; CCIX 互联规范与一致性总线）。
    \end{itemize}
    \item \textbf{领域专用架构}：
    \begin{itemize}
        \item 机器学习专用硬件加速器（脉动阵列 Systolic Array、数据流数据通路 Dataflow Datapath）。
        \item 网络内计算：可编程交换芯片（P4 架构）与智能网卡。
        \item 面向新兴前沿应用的低功耗专用嵌入式系统。
        \item 类脑/神经拟态计算架构（Neuromorphic Computing，脉冲神经网络 SNN、存算一体交叉阵列与忆阻器芯片）。
    \end{itemize}
    \item \textbf{先进半导体封装与异构集成方案}：三维集成电路（3DIC）、芯粒技术、硅中介层与通用芯粒互连标准。
\end{itemize}

\subsubsection{示范性学习成果与能力目标}
\begin{enumerate}
    \item \textbf{[评估]} 跟踪并研读包含典型异构并行微架构（如 CPU-SIMD、GPU 流多处理器 SM、NPU 脉动阵列）的现代异构系统集成框图，精准标注其数据流动与控制模型的核心差异。
    \item \textbf{[解释]} 深刻阐述在多处理器与多加速器异构系统中独有而在单处理器单核系统中不存在的存储管理难题（如非一致内存访问 NUMA 效应、跨设备虚拟地址翻译、异构缓存一致性风暴），并说明软硬件协同应对策略。
    \item \textbf{[评估]} 深入论证系统背板总线、处理器-内存直连通道以及基于高速网络的远程直接内存访问（RDMA）在访存延迟与吞吐量上的阶跃差异及其对分布式并行程序性能的直接制约。
    \item \textbf{[评估]} 针对给定的实际算法负载特征（如计算密度比 Arithmetic Intensity、访存模式规则性与控制流分支发散度），科学论证在何种场景下应当引入领域专用加速器（DSA）而非依赖通用多核 CPU。
    \item \textbf{[解释]} 系统剖析并阐述向量处理单元与标量处理单元在寄存器组织、指令流水与内存 Gather/Scatter 机制上的体系结构设计哲学差异。
    \item \textbf{[评估]} 全面评估存算一体（PIM）体系结构的工程优缺点：系统权衡其在粉碎“冯·诺依曼访存墙”、降低数据搬运能耗等方面的巨大收益，以及在模拟计算噪声、热耗散集中、制造良率与现有编译软件栈兼容性方面面临的技术挑战。
\end{enumerate}

\subsection{量子体系结构}
\textbf{【学时建议】}：3 领域核心学时。

本单元面向未来颠覆性计算前沿，构建量子物理基础与量子微架构体系的认知桥梁：

\subsubsection{领域核心主题}
\begin{itemize}
    \item \textbf{量子计算物理基础原理}：
    \begin{itemize}
        \item 微观粒子的波粒二象性原理。
        \item 双缝干涉实验中的海森堡测不准/不确定性原理。
        \item 什么是量子比特？量子态叠加与测量坍缩；以单光子作为物理量子比特的实现。
        \item 双量子比特系统：量子纠缠、贝尔基态与无超光速通信定理。
        \item 量子力学公理系统：态矢量叠加公理、投影测量公理、闭系统酉演化变换。
        \item 量子线路模型的单量子比特逻辑门：Pauli-X 门（非门）、Pauli-Z 门（相位翻转门）、Hadamard 门（H 门，创建叠加态）。
        \item 双量子比特受控门（如 CNOT 门）与张量积运算；量子门矩阵代数表示。
        \item 量子不可克隆定理；量子隐形传态协议。
    \end{itemize}
    \item \textbf{经典与基础量子算法}：
    \begin{itemize}
        \item 基础预言机算法：Bernstein-Vazirani 算法、Simon 周期查找算法。
        \item 利用马赫-曾德尔干涉仪光学实现 Deutsch-Jozsa 算法。
        \item 量子质因数分解算法（Shor's Algorithm）核心流程。
        \item 量子无序数据库加速搜索算法（Grover's Algorithm）。
    \end{itemize}
    \item \textbf{量子硬件实现与微架构工程}：
    \begin{itemize}
        \item 量子比特的主流物理实现路线（当前主流九大物理实现载体：超导量子电路、捕获离子阱、半导体量子点硅自旋、光量子芯片、中性冷原子阵列、拓扑准粒子、金刚石氮-空位色心 NV、液氦表面电子等）。
        \item 量子处理单元（QPU）的经典微电子控制架构（低温同轴线缆、微波脉冲生成、高精度 DAC/ADC 控制逻辑与实时反馈回路）。
        \item 量子误差缓解与量子纠错：表面码、含噪声中等规模量子（NISQ）时代的现实局限与容错通用量子计算远景。
    \end{itemize}
    \item \textbf{新兴前沿应用方向}：
    \begin{itemize}
        \item 抗量子后量子密码学（Post-quantum cryptography / PQC，如基于格密码的算法过渡）。
        \item 量子互联网（The Quantum Internet：量子中继器、长距离纠缠分发与量子密钥分发 QKD）。
        \item 绝热量子计算 (AQC)与量子退火。
    \end{itemize}
\end{itemize}

\subsubsection{示范性学习成果与能力目标}
\begin{enumerate}
    \item \textbf{[解释]} 深入阐释微观量子物体的波粒二象性运动图景：以粒子形态产生，以波函数形式在空间中传播演化，并在测量探测时以遵从波函数模平方概率分布的粒子形态呈现。
    \item \textbf{[解释]} 借助实验时序图，深入评述微观量子系统的内禀概率本质与观测介入导致的系统态坍缩特性。
    \item \textbf{[评估]} 严谨表述对量子纠缠的科学认知：纠缠可用于建立非定域的强非经典关联，但依据“无信号定理”，物理上绝对无法利用量子纠缠进行超光速的有效信息通信。
    \item \textbf{[解释]} 深刻阐明量子并行性的运行机理，并针对量子测量的概率特性，分析相长干涉与相消干涉在筛选与放大正确答案概率中的核心设计机制。
    \item \textbf{[应用]} 在给定的 Shor 算法微代码或线路图上，准确标注量子傅里叶采样（QFT）在提取量子叠加态相位与模周期中的关键运算角色。
    \item \textbf{[开发]} 在开源量子计算仿真平台（如 Qiskit / Cirq）中编写 Shor 算法仿真程序，并撰写工程报告，深入辨析经典计算模块（数论模幂规约、连分数展开）与量子加速核心子模块在算法整体中的协同边界。
    \item \textbf{[解释]} 系统梳理并详细阐明当前各主流量子比特实现路线（如超导、离子阱、硅自旋、光量子、中性原子、拓扑等）的物理工作机制、相干保真度指标与可扩展微加工壁垒。
    \item \textbf{[评估]} 深入辨析绝热量子计算（AQC）/量子退火与基于逻辑门的通用线路模型量子计算在硬件体系结构、控制机制与适用计算问题类别（如组合优化与退火基态搜索 vs 通用代数运算）上的本质差异。
    \item \textbf{[评估]} 深刻论证与评述以下关键体系结构命题：“在现代异构计算系统的宏观视角下，量子处理单元（QPU）本质上是一种由经典宿主计算机通过特定总线进行控制与卸载计算的异构专用协处理器微架构（类似于 FPGA 或 GPU）”。
\end{enumerate}

\section{专业素养与思维倾向}

在计算机体系结构与组织的研习与工程实践中，重点塑造学生以下维度的专业心智素养与职业特质：
\begin{itemize}
    \item \textbf{自主求知与深度钻研}：具备日益增强的自主求知欲与主动性，能自觉主动研读繁复的芯片体系结构白皮书、底层汇编规范、微架构参考手册以及工业级硬件设计文档，持续拓宽软硬件交叉知识边界。
    \item \textbf{前瞻洞察与跨界融通}：展现积极主动且独立探索的工程素养，能跨越传统学科界限，无缝整合体系结构、编译优化、操作系统调度与算法分析等多领域知识，精准洞察与驾驭计算机系统的底层运行逻辑。
    \item \textbf{突破创新的架构思维}：勇于超越平凡直觉与常规解决方案，面对苛刻的硬件算力、带宽瓶颈与功耗墙约束，能因地制宜、创造性地发掘并利用底层硬件专有微架构特性（如 SIMD 向量指令、片上共享存储、异步 DMA 流水等）以达成极致能效。
    \item \textbf{严谨敬业与工程伦理操守}：秉持高度的职业审慎态度与道德责任感。现代计算机系统（特别是直接贴附或植入人体的生物传感器芯片，如实时血糖监测与心脏起搏设备）与人身安全直接相连，工程师必须在硬件芯片与微架构底层将用户隐私、物理安全性与系统可靠性置于至高无上的优先级。
\end{itemize}

\section{数学基础要求}

计算机体系结构领域对离散数学、线性代数及统计物理有着多层次的数学依赖：

\subsection{必须掌握的核心数学}
\begin{itemize}
    \item \textbf{离散数学}：集合论、二元关系、布尔逻辑与谓词运算（用于数字逻辑综合与状态机设计）、初等数论（同余定理、原根与模算术，用于计算机算术、差错控制编码与硬件密码加速）。
    \item \textbf{线性代数}：算术基础运算、矩阵乘法与多维张量操作（用于 GPU/TPU 脉动阵列加速器微架构设计）、矩阵特征值与酉算符分解（用于量子逻辑门矩阵代数）。
\end{itemize}

\subsection{进阶推荐数学背景}
\begin{itemize}
    \item \textbf{面向量子计算的数理基础}：基础概率论与随机变量、三角函数、复数向量空间与复数内积运算、欧拉公式（$e^{i\theta} = \cos\theta + i\sin\theta$）。
    \item \textbf{系统性能评估与试验设计方法}：数理统计基础、正态分布与置信区间、析因实验设计与方差分析（用于多微架构参数调优、性能瓶颈敏感度消融实验）。
\end{itemize}

\section{交叉主题与跨学科共享概念}

\subsection{跨学科共享概念}
\begin{itemize}
    \item \textbf{虚拟内存}：体系结构/层次存储系统 与 操作系统/内存管理（操作系统/内存管理）深度共享。
    \item \textbf{差错控制编码、数据压缩与数据完整性}：体系结构/层次存储系统 与 系统基础/Software Reliability through Redundancy（系统基础/基于冗余的可靠性）深度共享。
    \item \textbf{通信互联与片上网络}：体系结构/接口与通信 与 网络通信/网络导论与体系结构（网络与通信导论）深度共享。
    \item \textbf{处理器微架构设计}：体系结构/异构架构 与 专用平台/移动平台（移动平台）及 专用平台/Industrial（工业平台）交叉。
    \item \textbf{领域专用加速架构}：体系结构/异构架构 与 人工智能/Advanced Machine Learning（高级机器学习）交叉。
    \item \textbf{新兴硬件用户接口}：体系结构/异构架构 与 人机交互/Human Factors \& Security（人机交互/人类因素与安全）交叉。
    \item \textbf{底层并行计算系统}：体系结构/汇编级机器组织 与 并行分布式/并行体系结构（并行计算/并行体系结构）交叉。
\end{itemize}

\subsection{跨学科横切主题}
\begin{itemize}
    \item \textbf{多层级并行性}：指令级并行（ILP）、数据级并行（SIMD/向量化）、线程级并行（TLP）与请求级并行（RLP）。
    \item \textbf{系统资源分配、硬件调度与空间隔离}。
    \item \textbf{性能与能源效率的永恒权衡}。
    \item \textbf{底层计算系统的社会与生态环境影响}：数据中心碳排放、高能耗芯片制冷挑战、暗硅现象与电子硬件生命周期管理。
\end{itemize}

\section{课程设计与开课规划建议}

联合工作组针对不同层次的培养目标推荐两套典型的开课规划方案：

\subsection{方案一：计算机体系结构入门课程 (建议约 42 学时)}
本方案面向本科通用计算机专业，构建完整的软硬件交界面认知：
\begin{itemize}
    \item 社会伦理/计算机发展简史 [2 学时]
    \item 体系结构/机器级数据表示 [2 学时]
    \item 体系结构/汇编级机器组织 [2 学时]
    \item 体系结构/层次存储系统 [10 学时]
    \item 操作系统/内存管理 [10 学时]
    \item 体系结构/接口与通信 [4 学时]
    \item 体系结构/异构架构 [5 学时]
    \item 并行分布式/程序与执行 [4 学时]
    \item 社会伦理/伦理分析方法 [3 学时]
\end{itemize}
\textbf{【先修要求】}：离散数学（集合、二元关系、逻辑运算、初等数论）。\\
\textbf{【能力目标】}：完成本课程的学生应能全面理解现代计算机系统的底层基础架构，深刻领会 Cache 高速缓存、流水线冲突以及硬件内存管理机制所面临的系统级挑战。

\subsection{方案二：高级计算机系统综合课程 (建议约 38 学时)}
本方案面向计算机系统方向的高年级拔尖创新人才与研究型培养：
\begin{itemize}
    \item 社会伦理/计算机发展简史 [2 学时]
    \item 社会伦理/隐私保护 [3 学时]
    \item 软件工程/软件架构与设计 [4 学时]
    \item 并行分布式/并行算法与应用 [4 学时]
    \item 体系结构/异构架构 [4 学时]
    \item 操作系统/角色与定位 [3 学时]
    \item 体系结构/层次存储系统 [7 学时]
    \item 体系结构/性能与能效 [5 学时]
    \item 网络通信/网络化应用与套接字编程 [5 学时]
\end{itemize}
\textbf{【先修要求】}：离散数学（集合、二元关系、逻辑运算、初等数论）。\\
\textbf{【能力目标】}：完成本课程的学生应能系统洞察现代计算硬件体系结构的核心设计哲学，具备在后摩尔时代驾驭硬件并行性与异构算力、持续提升系统性能与能效的高维系统工程素养。

\section{真实能力任务规范}

\begin{taskbox}[任务 1：应用程序中的高速缓存性能瓶颈评估与优化]
\textbf{【任务情境】}：某企业核心高频业务系统的延迟指标出现严重长尾波动，开发团队怀疑是由底层的 Cache 缺失与颠簸导致。\\
\textbf{【能力陈述】}：批判性地剖析应用程序在高速缓存层面的访存性能瓶颈（命中率、行冲突、假共享），运用微架构性能计数器（PMU）进行量化剖析，并撰写总结核心瓶颈成因与重构建议的权威评估报告。\\
\textbf{【能力领域】}：系统 / 应用 \quad \textbf{【能力单元】}：系统开发 / 部署 / 评估 / 改进\\
\textbf{【关联知识单元】}：体系结构/层次存储系统、体系结构/性能与能效、操作系统/内存管理\\
\textbf{【达标等级】}：\textbf{评估}
\end{taskbox}

\begin{taskbox}[任务 2：边缘嵌入式设备上机器学习模型部署的“每瓦性能比”评估]
\textbf{【任务情境】}：在电池供电的低功耗移动机器人上部署云端训练完成的目标检测神经网络模型，需在帧率与续航之间达成最优平衡。\\
\textbf{【能力陈述】}：在目标嵌入式硬件设备上实测并评估运行轻量化机器学习模型的端到端推断吞吐量、芯片动态功耗与热耗散表现，计算其“每瓦性能比”（Performance-per-Watt），指导模型量化与微架构剪枝策略。\\
\textbf{【能力领域】}：系统 / 应用 \quad \textbf{【能力单元】}：部署 / 评估 / 改进\\
\textbf{【关联知识单元】}：体系结构/性能与能效、专用平台/移动平台\\
\textbf{【达标等级】}：\textbf{评估与开发}
\end{taskbox}

\begin{taskbox}[任务 3：将 CPU 应用程序重构移植至专用硬件加速器]
\textbf{【任务情境】}：传统在 x86 CPU 上运行的大规模科学模拟或图像处理程序算力不足，需将其高性能关键计算循环迁移至现代硬件加速器上。\\
\textbf{【能力陈述】}：综合运用系统架构设计知识，将通用 CPU 应用程序的核心密集运算代码重构并移植至专用硬件加速器（GPU、TPU 或 NPU），合理规划设备端内存分配与主机-设备数据传输流水线，并精确评估计算加速比。\\
\textbf{【能力领域】}：软件 / 系统 \quad \textbf{【能力单元】}：需求 / 设计 / 开发 / 测试 / 技术文档\\
\textbf{【关联知识单元】}：体系结构/异构架构、并行分布式/并行体系结构、系统基础/系统设计\\
\textbf{【达标等级】}：\textbf{评估与开发}
\end{taskbox}

\begin{taskbox}[任务 4：构建性能基准测试工具以诊断并消除程序 I/O 吞吐瓶颈]
\textbf{【任务情境】}：大规模数据处理程序在吞吐量扩展时遭遇未知物理阻滞，亟需确定系统瓶颈究其在于 CPU 计算能力还是外部 I/O 通信通路。\\
\textbf{【能力陈述】}：开发或熟练配置一套系统级性能基准测试与性能剖析工具，实时监测磁盘总线、网络套接字与内存 DMA 数据流，量化识别 I/O 阻塞节点，并评估通过消除 I/O 瓶颈所获得的系统级性能加速收益。\\
\textbf{【能力领域】}：软件 / 系统 \quad \textbf{【能力单元】}：开发 / 评估 / 系统维护\\
\textbf{【关联知识单元】}：体系结构/接口与通信、体系结构/性能与能效、系统基础/资源分配与调度、系统基础/系统性能\\
\textbf{【达标等级】}：\textbf{评估与解释}
\end{taskbox}

\begin{taskbox}[任务 5：评估 CPU-GPU 异构统一内存管理中的缺页异常及其性能损耗]
\textbf{【任务情境】}：在现代异构统一虚拟内存架构中，过多的主机与加速器间隐式页面迁移导致应用程序性能剧烈退化。\\
\textbf{【能力陈述】}：运用操作系统内存管理与计算机体系结构原理，深入剖析 CPU 与 GPU 之间统一虚拟内存的缺页异常触发机制、页面迁移开销与 PCIe/NVLink 总线拥堵效应，并提出显式预取与提示优化方案。\\
\textbf{【能力领域】}：系统 \quad \textbf{【能力单元】}：系统设计 / 开发 / 评估 / 改进\\
\textbf{【关联知识单元】}：体系结构/层次存储系统、体系结构/性能与能效、体系结构/异构架构、操作系统/内存管理\\
\textbf{【达标等级】}：\textbf{应用与评估}
\end{taskbox}

\begin{taskbox}[任务 6：连续体征监测医疗设备在硬件与系统级的数据隐私与安全影响评估]
\textbf{【任务情境】}：研发一款贴附于患者身体、具备高频连续生理体征数据采集能力的穿戴式医疗设备，需规避硬件侧信道泄露与隐私窃取。\\
\textbf{【能力陈述】}：应用底层计算系统的专业技术知识，权衡嵌入式传感芯片的数据采集精度、片上数据最小化处理策略与硬件安全信任根机制，撰写关于医疗设备硬件与固件级数据隐私合规的技术白皮书。\\
\textbf{【能力领域】}：系统 / 应用 \quad \textbf{【能力单元】}：需求分析 / 系统设计 / 技术文档 / 社会议题调适\\
\textbf{【关联知识单元】}：体系结构/异构架构、社会伦理/隐私保护、编程语言/嵌入式计算与硬件接口\\
\textbf{【达标等级】}：\textbf{评估}
\end{taskbox}

\begin{taskbox}[任务 7：传感器硬件集成的软硬件协同驱动接口与模块设计]
\textbf{【任务情境】}：为新型高精度工业物联网传感硬件开发低延迟、高可靠的板级支持包（BSP）与操作系统驱动框架。\\
\textbf{【能力陈述】}：设计并实现面向传感器芯片的高层软件抽象 API 以及底层基于中断与总线寄存器的硬件交互驱动模块，达成无缝的软硬件系统集成与容错处理。\\
\textbf{【能力领域】}：系统 \quad \textbf{【能力单元】}：需求工程 / 架构设计 / 模块开发\\
\textbf{【关联知识单元】}：体系结构/接口与通信、体系结构/异构架构、操作系统/文件系统 API 与底层实现、软件开发基础/基本编程概念、软件工程/API 设计原则、专用平台/嵌入式系统\\
\textbf{【达标等级】}：\textbf{开发与评估}
\end{taskbox}

\begin{taskbox}[任务 8：系统软硬件架构选型论证与技术方案设计规范撰写]
\textbf{【任务情境】}：面向高可靠严苛工控领域，企业拟研发下一代综合计算平台，需全面论证处理器微架构选型与系统拓扑方案。\\
\textbf{【能力陈述】}：全面收集梳理复杂应用场景的计算需求、功耗预算与散热条件，权衡通用处理器与专用加速芯片配置，提出系统软硬件体系结构综合设计方案，并撰写高质量的技术决策规范文档。\\
\textbf{【能力领域】}：系统 / 应用 \quad \textbf{【能力单元】}：需求工程 / 架构设计 / 技术文档\\
\textbf{【关联知识单元】}：体系结构/机器级数据表示、体系结构/性能与能效、体系结构/异构架构、社会伦理/专业沟通与表达、软件工程/需求工程与产品规约\\
\textbf{【达标等级】}：\textbf{评估与解释}
\end{taskbox}

\begin{taskbox}[任务 9：自动化制造单元中高精度感知-执行机械臂系统的跨学科协同研发]
\textbf{【任务情境】}：在现代智能制造柔性装配线中，跨学科研发团队正在攻坚具备毫秒级视觉反馈与高精度伺服控制的智能协作机械臂系统。\\
\textbf{【能力陈述】}：在多学科复合型研发团队中高效沟通协作，负责机械臂底层嵌入式微控制器、总线通信协议与实时加速芯片的软硬件联合设计、集成调试与系统性能验证。\\
\textbf{【能力领域】}：软件 / 系统 / 应用 \quad \textbf{【能力单元】}：需求 / 设计 / 开发 / 测试 / 集成 / 评估\\
\textbf{【关联知识单元】}：体系结构/异构架构、专用平台/机器人平台、软件工程/团队协作\\
\textbf{【达标等级】}：\textbf{应用与开发}
\end{taskbox}

\begin{taskbox}[任务 10：深度利用底层多核 CPU 与层次存储系统的高性能并发程序开发]
\textbf{【任务情境】}：针对金融量化高频交易或超大规模网格计算，需彻底压榨现代多核多级缓存服务器硬件的每一周期算力。\\
\textbf{【能力陈述】}：透彻理解现代 CPU 核心流水线、乱序执行、SIMD 向量扩展与层次内存架构，设计并编写底层高性能并发程序，消除伪共享、优化缓存局部性、精细化调度多核线程以达成极致吞吐。\\
\textbf{【能力领域】}：系统 \quad \textbf{【能力单元】}：需求 / 设计 / 开发 / 性能调优\\
\textbf{【关联知识单元】}：体系结构/汇编级机器组织、体系结构/层次存储系统\\
\textbf{【达标等级】}：\textbf{评估与开发}
\end{taskbox}

\section{指导委员会成员名单}

\begin{itemize}
    \item \textbf{主席}：Marcelo Pias，巴西里奥格兰德联邦大学, Brazil
    \item \textbf{工作组成员}：
    \begin{itemize}
        \item Brett A. Becker，爱尔兰都柏林大学学院, Ireland
        \item Mohamed Zahran，美国纽约大学, USA
        \item Monica D. Anderson，美国阿拉巴马大学, USA
        \item 向乔，厦门大学
        \item Adrian German，美国印第安纳大学布卢明顿分校, USA
    \end{itemize}
\end{itemize}

\section{附录：核心主题与技能要求等级对照表}

\begin{longtable}{p{0.18\textwidth} p{0.48\textwidth} c c c}
\toprule
\textbf{知识单元} & \textbf{核心知识点与具体主题} & \textbf{认知等级} & \textbf{分类} & \textbf{学时} \\
\midrule
\endfirsthead
\toprule
\textbf{知识单元} & \textbf{核心知识点与具体主题} & \textbf{认知等级} & \textbf{分类} & \textbf{学时} \\
\midrule
\endhead
\bottomrule
\endfoot
\multirow{2}{*}{\makecell[l]{体系结构/数字逻辑与\\数字系统}} 
& 体系结构演进历史；组合逻辑与时序逻辑基本构建块；硬件与软件多层抽象体系结构 & 解释 & 领域核心 & \multirow{2}{*}{3} \\
\cmidrule{2-4}
& 硬件 CAD 工具链；高层次综合（RTL、Verilog/VHDL/Chisel）；SoC 设计流程；物理与电气约束（延迟、扇入扇出、功耗、光速） & 评估 & 领域核心 & \\
\midrule
\makecell[l]{体系结构/机器级\\数据表示} 
& 位、字节与字；定点数与浮点数（IEEE 754）；原码与补码运算及溢出；非数值字符编码；记录/结构体对齐与多维数组内存排布 & 应用 & 计算机核心 & 1 \\
\midrule
\multirow{2}{*}{\makecell[l]{体系结构/汇编级\\机器组织}} 
& 经典冯·诺依曼架构；控制单元与指令周期；SIMD 与 MIMD 导论及 Flynn 分类；共享内存多核多处理器组织 & 解释 & 计算机核心 & 1 \\
\cmidrule{2-5}
& 指令集体系结构（ISA：x86/ARM/RISC-V）；寻址模式；机器码与汇编程序设计；子程序调用机制与栈帧；I/O 与中断；内存段划分 & 开发 & 领域核心 & 2 \\
\midrule
\multirow{2}{*}{\makecell[l]{体系结构/层次存储系统}} 
& 时间与空间局部性；主存组织（DRAM/DDR）；持久化存储（SSD/HDD）；访存延迟、周期、带宽与交叉存取；虚拟内存支持；硬件可靠性与纠错码；存算一体 & 解释 & 计算机核心 & \multirow{2}{*}{6} \\
\cmidrule{2-4}
& 高速缓存 Cache（地址映射、块大小、替换与写策略）；多处理器缓存一致性（监听/目录协议、MESI/MOESI 状态机） & 评估 & 计算机核心 & \\
\midrule
\makecell[l]{体系结构/接口与通信} 
& I/O 基础原理（握手、缓冲、轮询、中断驱动）；向量与优先级中断结构；外存储物理组织；系统总线协议、仲裁与 DMA；片上网络 & 解释 & 计算机核心 & 1 \\
\midrule
\multirow{2}{*}{\makecell[l]{体系结构/功能部件组织}} 
& 数据通路设计；控制单元实现（硬布线逻辑与微程序控制器）；五级指令流水线实现；流水线冒险（结构、数据、控制冒险）检测与消解 & 开发 & 领域核心 & \multirow{2}{*}{2} \\
\cmidrule{2-4}
& 深入指令流水线（时空图）；指令级并行（ILP）核心概念与超标量架构基础 & 解释 & 领域核心 & \\
\midrule
\multirow{2}{*}{\makecell[l]{体系结构/性能与能效}} 
& 性能与能耗综合评估导论；分支预测技术；推测执行；乱序执行与 Tomasulo 算法；硬件预取技术 & 评估 & 领域核心 & 2 \\
\cmidrule{2-5}
& 向量处理器与 GPU 增强；多线程硬件支持（硬件锁、同步与屏障）；可扩展性定律；VLIW 替代架构；DVFS 动态能效调节；暗硅挑战 & 解释 & 领域核心 & 1 \\
\midrule
\makecell[l]{体系结构/异构架构} 
& SIMD/MIMD 异构架构；异构存储（共享 vs 分布式、CXL 缓存一致性）；领域专用架构 DSA（机器学习加速器、网络内计算、类脑计算）；先进封装（3DIC 与 Chiplets 芯粒） & 解释 & 领域核心 & 3 \\
\midrule
\makecell[l]{体系结构/量子体系结构} 
& 量子物理原理（波粒二象性、测不准原理、Qubit、叠加态、纠缠、公理体系、单/双量子比特门、不可克隆定理）；量子算法（Bernstein-Vazirani、Simon、Shor、Grover）；硬件实现（九大物理路线、经典控制、NISQ）；新兴应用（抗量子密码 PQC、量子互联网、量子退火） & 解释 & 领域核心 & 3 \\
\bottomrule
\end{longtable}
